\documentclass[12pt]{article}
\usepackage[utf8]{inputenc}
\usepackage[T1]{fontenc}
\usepackage{hyperref}
\title{Veddah Language Resources for K-12 Curriculum Development}
\author{Compiled by Claude Sonnet 5 (claude-sonnet-5) for Andr\'as Kornai}
\date{}
\begin{document}
\maketitle
\section{Introduction}
Veddah (ISO 639-3: ved) is an endangered language of Sri Lanka. Search results describe it as classed by UNESCO as definitely endangered; one source (Verbix) lists it as dormant with no speakers. Linguists have debated whether it is a Sinhala dialect, an independent language, or a creole. Very little documentation was found.
\section{Primary Educational Resources}
No primers, textbooks or school programmes for Veddah were found in search results.
\section{Linguistic and Reference Materials}
\begin{itemize}
\item Wikipedia, Vedda language (overview of status, history of study): \url{https://en.wikipedia.org/wiki/Vedda_language}
\item Verbix Endangered Languages page for Veddah: \url{https://docs.verbix.com/EndangeredLanguages/Veddah}
\item Glottolog reference entry (Dharmadasa 1990, ethnographic rather than linguistic): \url{https://glottolog.org/resource/reference/id/313723}
\end{itemize}
\section{Religious and Community Resources}
None found.
\section{Digital and Audio Resources}
No audio or digital corpora found.
\section{Curriculum Development Context}
Since no dictionary, grammar or teaching material was found, any curriculum would depend on new documentation work with Vedda community members and on specialist literature not located here (for example the 1959 field study mentioned in the sources).
\end{document}